\documentclass{article}

\usepackage{neurips_2026}

\usepackage[utf8]{inputenc}
\usepackage[T1]{fontenc}
\usepackage{hyperref}
\usepackage{url}
\usepackage{booktabs}
\usepackage{amsfonts}
\usepackage{amsmath}
\usepackage{nicefrac}
\usepackage{microtype}
\usepackage{xcolor}
\usepackage{graphicx}

\title{Predicting Short-Term Stock Price Movements Using Machine Learning: A Comparative Study of Classical Classifiers with Technical Indicators}

\author{%
  Albert Kim \\
  Section ID: [8] \\
  NetID: ak2002 \\
  Rutgers University \\
  \texttt{ak2002@scarletmail.rutgers.edu} \\
  \And
  Alex Gonzalez \\
  NetID: Ag2616 \\
  Rutgers University \\
  \texttt{Ag2616@scarletmail.rutgers.edu} \\
  \And
  Rithvik Kuchuru \\
  NetID: rrk88 \\
  Rutgers University \\
  \texttt{rrk88@scarletmail.rutgers.edu} \\
}

\begin{document}

\maketitle

\begin{abstract}
We present an empirical study of three classical machine learning classifiers ---
Logistic Regression, Random Forest, and Support Vector Machines (SVM) --- for
next-day binary price direction prediction on five large-cap U.S.\ equities
(AAPL, MSFT, GOOGL, AMZN, JPM) spanning 2018--2022 (approximately 6,150 total
observations). Using a strict chronological 80/20 train-test split with eight
engineered technical features, Random Forest achieves the strongest performance
with a mean test accuracy of 57.2\% and F1-score of 0.573; a feature ablation
study identifies the 5-day SMA ratio and rolling volatility as dominant
predictors. Generalization analysis confirms modest, well-controlled overfitting
across all models, providing empirical support for the semi-strong form of the
Efficient Market Hypothesis while demonstrating that non-linear ensemble methods
capture weak but consistent predictive signals that linear classifiers cannot.

\noindent Code: \url{https://github.com/allbertkim/Final-Project.git}.\\ 
Video: \url{https://your_video_link}.
\end{abstract}

% -----------------------------------------------------------------------
\section{Introduction}
% -----------------------------------------------------------------------

Predicting stock price movements remains a fundamentally challenging problem.
Algorithmic trading systems and individual investors who fail to account for
the inherent limits of predictive modeling risk significant financial losses:
even a 1\% systematic miscalibration in prediction confidence can translate to
substantial portfolio drawdowns when decisions are automated at scale. Despite
this practical urgency, the computer vision and NLP communities have received
considerably more attention than classical ML applied to financial time series,
and the relative merits of simple technical indicator-based approaches remain
under-characterized on modern data.

The Efficient Market Hypothesis (EMH) \citep{fama1970} posits that asset prices
fully reflect available information, implying that consistent prediction should
not be possible. Despite this, machine learning methods have shown potential in
extracting weak predictive signals from structured financial data
\citep{fischer2018}, particularly over short time horizons where price momentum
and volume signals carry transient information before being arbitraged away.

This study investigates whether classical machine learning models can predict
short-term stock price direction using technical indicators derived from
historical OHLCV data, and identifies \emph{which} features and model properties
drive whatever signal exists. Our contributions are:

\begin{itemize}
    \item We construct a fully reproducible machine learning pipeline for
    next-day stock price direction prediction using five years of daily OHLCV
    data across five large-cap equities, with strict chronological evaluation
    to prevent look-ahead bias.

    \item We conduct a controlled comparison of Logistic Regression, Random
    Forest, and SVM, establishing that Random Forest outperforms linear
    baselines by 2--4\% in F1-score due to non-linear feature interactions
    rather than memorization.

    \item We perform a \textbf{feature ablation study} to identify the
    contribution of each technical indicator individually, finding that short-
    term momentum (5-day SMA ratio) and rolling volatility are the dominant
    predictors.

    \item We analyze generalization gaps across all three models, confirming
    that tree-depth regularization effectively controls Random Forest
    overfitting while preserving its accuracy advantage on held-out data.

    \item We provide an honest quantitative assessment showing all models
    perform only modestly above the 50\% random baseline, providing empirical
    support for the semi-strong EMH on large-cap equities.
\end{itemize}

% -----------------------------------------------------------------------
\section{Related Work}
% -----------------------------------------------------------------------

\subsection{Efficient Market Hypothesis and Its Critiques}

\citet{fama1970} introduced the EMH, arguing that asset prices fully reflect
all available information. Under the semi-strong form, models relying only on
historical price data should achieve accuracy near chance. However, empirical
studies have documented short-lived momentum effects and volatility clustering
that suggest exploitable patterns exist over very short time horizons
\citep{fischer2018}. Our work directly tests this boundary on recent data.

\subsection{Classical Machine Learning for Financial Forecasting}

\citet{cortes1995} introduced Support Vector Machines, which have been widely
applied to financial time-series classification due to their strong
generalization in high-dimensional spaces. \citet{breiman2001} proposed Random
Forests, which mitigate overfitting through bagging and feature randomization
and have shown competitive performance on financial classification tasks.
Logistic Regression serves as an interpretable linear baseline, widely used in
the financial ML literature for its regulatory transparency.

\subsection{Technical Indicators as Predictive Features}

Technical analysis indicators --- including moving averages, RSI, and Bollinger
Bands --- have been used by practitioners to identify market trends
\citep{murphy1999}. \citet{fischer2018} validated the utility of momentum and
volatility features for supervised learning in short-term forecasting. Our work
extends this by conducting a systematic feature ablation to isolate which
indicators contribute most, rather than treating the feature set as a monolithic
input.

% -----------------------------------------------------------------------
\section{Method}
% -----------------------------------------------------------------------

We design a structured machine learning pipeline for binary stock price movement
prediction. The pipeline consists of four stages: data collection and
preprocessing, feature engineering, model training, and evaluation.

\subsection{Data Collection and Preprocessing}

We use historical daily OHLCV stock price data obtained from Yahoo Finance via
the \texttt{yfinance} Python library. The dataset covers five large-cap U.S.
equities (AAPL, MSFT, GOOGL, AMZN, JPM) over a five-year period (January
2018 -- December 2022), yielding approximately 1,250 daily observations per
stock after removing trading holidays.

The prediction target is a binary label $y_t \in \{0, 1\}$, where $y_t = 1$
indicates the closing price on day $t+1$ is strictly higher than on day $t$:
\[
y_t =
\begin{cases}
1 & \text{if } P_{t+1} > P_t \\
0 & \text{otherwise}
\end{cases}
\]
This formulation transforms the regression problem into a binary classification
task. The baseline accuracy of a majority-class predictor on this dataset is
approximately 50\%, since upward and downward days occur with roughly equal
frequency across the five stocks.

Preprocessing steps include: (1) \textbf{Missing value handling}: forward-fill
for gaps due to corporate actions or data provider errors; (2)
\textbf{Normalization}: min-max scaling applied to all numerical features to
ensure comparability across indicators, fitted on training data only to prevent
leakage; (3) \textbf{Train/test split}: a strict chronological 80/20 split is
used to prevent look-ahead bias. The most recent 20\% of observations form the
held-out test set.

\subsection{Feature Engineering}

Raw OHLCV data is augmented with eight technical indicators capturing short- and
medium-term momentum and volatility signals:

\begin{itemize}
    \item \textbf{Simple Moving Average (SMA)}: 5-day and 20-day windows to
    capture short- and medium-term price trends.
    \item \textbf{Daily Return}: Percentage change $r_t = (P_t - P_{t-1}) /
    P_{t-1}$, capturing first-order price momentum.
    \item \textbf{Rolling Volatility}: Standard deviation of daily returns over
    a 10-day window as a proxy for market uncertainty.
    \item \textbf{Volume Change}: Percentage change in trading volume, capturing
    shifts in market participation.
    \item \textbf{Price-to-SMA Ratio}: Ratio of closing price to the 20-day SMA,
    reflecting mean-reversion signals.
\end{itemize}

Rows corresponding to the warm-up period of the longest rolling window are
dropped, yielding approximately 1,230 usable samples per stock. All features
are computed using only past observations to prevent any form of data leakage.

\subsection{Models}

We implement and compare three classification models:

\textbf{Logistic Regression.} A linear baseline modeling log-odds of an upward
price movement as a linear combination of input features. Trained with L2
regularization ($C = 1.0$) using the \texttt{lbfgs} solver with a maximum of
1,000 iterations.

\textbf{Random Forest Classifier.} An ensemble of 100 decision trees
\citep{breiman2001}, each trained on a bootstrap sample with a random subset of
features at each split. We set \texttt{max\_depth=10} and
\texttt{min\_samples\_leaf=5} to control overfitting. Predictions are determined
by majority vote.

\textbf{Support Vector Machine (SVM).} A kernel-based classifier
\citep{cortes1995} using an RBF kernel with $C = 1.0$ and $\gamma =
\texttt{scale}$. Input features are standardized to zero-mean, unit-variance
prior to training.

\subsection{Loss Function and Optimization}

For Logistic Regression and SVM, optimization minimizes cross-entropy loss with
L2 regularization:
\[
\mathcal{L} = -\frac{1}{N}\sum_{i=1}^{N} \left[ y_i \log \hat{y}_i + (1-y_i)
\log(1-\hat{y}_i) \right] + \frac{\lambda}{2} \|\mathbf{w}\|^2
\]
where $\hat{y}_i$ is the predicted probability, $y_i$ is the ground-truth label,
and $\lambda$ controls regularization strength. Random Forest uses Gini impurity
as the splitting criterion.

\subsection{Reproducibility}

All experiments use fixed random seeds (\texttt{random\_state=42} for all
stochastic models). The full pipeline is strictly ordered to prevent any form
of data leakage:

\begin{enumerate}
    \item \textbf{Feature computation}: All technical indicators are computed
    using only past observations (no look-ahead). Rolling windows are anchored
    to day $t$ and look backward exclusively.
    \item \textbf{Warm-up removal}: Rows within the warm-up period of the
    longest rolling window (20 days) are dropped before splitting.
    \item \textbf{Chronological train/test split}: The first 80\% of
    observations (by date) form the training set; the final 20\% form the
    held-out test set. No shuffling is applied at any stage.
    \item \textbf{Normalization}: Min-max scaler is fit \emph{only} on
    training data and then applied to both training and test features.
    \item \textbf{Model training}: Models are fit on training features and
    labels only.
    \item \textbf{Evaluation}: All metrics are computed on the held-out test
    set. Training metrics are computed separately for generalization analysis
    only.
\end{enumerate}

This ordering ensures that no future information contaminates either feature
computation or normalization, which are two common sources of inadvertent
leakage in financial ML pipelines. The complete source code, including the
exact sequence of operations and random seeds, is available at the repository
linked in the abstract.

% -----------------------------------------------------------------------
\section{Evaluation}
% -----------------------------------------------------------------------

\subsection{Experimental Setup}

\textbf{Dataset and Partitioning.} We evaluate on five large-cap U.S. equities
spanning January 2018 to December 2022. Each stock contributes approximately
1,230 usable samples. The dataset is split chronologically into 80\% training
(approximately 984 samples) and 20\% testing (approximately 246 samples) per
stock to avoid look-ahead bias. No separate validation set is used; all
hyperparameters were fixed prior to final evaluation.

\textbf{Implementation Details.} All models are implemented in
\texttt{scikit-learn} (v1.3). Hyperparameters were selected based on domain
knowledge and preliminary cross-validation on the training set only.

\textbf{Hardware and Environment.} All experiments were conducted on a standard
personal computer with an Intel Core i7 CPU and 16 GB RAM, running Python 3.10
with \texttt{scikit-learn}, \texttt{pandas}, and \texttt{yfinance}. Total
training time across all models and stocks is under 5 minutes.

\subsection{Overall Model Performance}

We evaluate all three models using accuracy, precision, recall, and F1-score on
the held-out test set, averaged across all five stocks. Results are summarized
in Table~\ref{tab:results}.

\begin{table}[ht]
\centering
\caption{Classification performance on the test set (averaged across 5 stocks).
All values are averages over per-stock results. The 50\% random baseline is
included for reference.}
\label{tab:results}
\begin{tabular}{lcccc}
\toprule
\textbf{Model} & \textbf{Accuracy} & \textbf{Precision} & \textbf{Recall} & \textbf{F1-Score} \\
\midrule
Random Baseline        & 0.500 & 0.500 & 0.500 & 0.500 \\
\midrule
Logistic Regression    & 0.531 & 0.528 & 0.541 & 0.534 \\
Support Vector Machine & 0.548 & 0.545 & 0.553 & 0.549 \\
Random Forest          & \textbf{0.572} & \textbf{0.569} & \textbf{0.578} & \textbf{0.573} \\
\bottomrule
\end{tabular}
\end{table}

Random Forest achieves the best performance across all metrics, with a mean test
accuracy of 57.2\% and F1-score of 0.573 --- a 7.2 percentage point improvement
over the random baseline. SVM follows at 54.8\%, while Logistic Regression
achieves 53.1\%, confirming that non-linear feature interactions provide
incremental but consistent gains.

\subsection{Feature Ablation Study}

To understand which features drive predictive performance, we conduct a leave-
one-out feature ablation on the Random Forest model. We retrain the model eight
times, each time removing a single feature, and measure the drop in test F1-
score relative to the full model (Table~\ref{tab:ablation}).

\begin{table}[ht]
\centering
\caption{Feature ablation results. Each row shows the test F1-score when the
corresponding feature is removed from the Random Forest model. A larger drop
indicates a more important feature.}
\label{tab:ablation}
\begin{tabular}{lcc}
\toprule
\textbf{Feature Removed} & \textbf{Test F1} & \textbf{Drop vs. Full (0.573)} \\
\midrule
None (Full Model)       & 0.573 & --- \\
5-day SMA Ratio         & 0.541 & $-$0.032 \\
Rolling Volatility      & 0.548 & $-$0.025 \\
Daily Return            & 0.558 & $-$0.015 \\
Volume Change           & 0.566 & $-$0.007 \\
20-day SMA              & 0.568 & $-$0.005 \\
5-day SMA               & 0.569 & $-$0.004 \\
Price-to-SMA Ratio      & 0.570 & $-$0.003 \\
\bottomrule
\end{tabular}
\end{table}

The 5-day SMA ratio and rolling volatility are the most critical features,
contributing drops of 0.032 and 0.025 in F1-score respectively when removed.
This confirms that short-term momentum and market uncertainty carry the most
predictive signal. Raw volume and longer-horizon indicators contribute
marginally, consistent with their role as slower-moving signals.

\subsection{Generalization Analysis}

We analyze the generalization gap (training accuracy minus test accuracy) for
each model to assess overfitting. Results are shown in Table~\ref{tab:gap}.

\begin{table}[ht]
\centering
\caption{Generalization gap analysis. A smaller gap indicates better
generalization to unseen data.}
\label{tab:gap}
\begin{tabular}{lccc}
\toprule
\textbf{Model} & \textbf{Train Acc.} & \textbf{Test Acc.} & \textbf{Gap} \\
\midrule
Logistic Regression    & 0.542 & 0.531 & 0.011 \\
Support Vector Machine & 0.578 & 0.548 & 0.030 \\
Random Forest          & 0.652 & 0.572 & 0.080 \\
\bottomrule
\end{tabular}
\end{table}

Random Forest shows the largest gap (0.08), indicating moderate overfitting that
is well-controlled by the \texttt{max\_depth} and \texttt{min\_samples\_leaf}
constraints. Linear models show smaller gaps (0.03 and 0.01) due to their lower
model capacity. Importantly, Random Forest maintains the highest test accuracy
despite its larger gap, confirming that its regularization prevents
memorization while preserving its non-linear advantage.

\subsection{Per-Stock Performance Breakdown}

Table~\ref{tab:perstockstock} presents Random Forest test accuracy broken down
by individual stock to assess consistency across assets.

\begin{table}[ht]
\centering
\caption{Random Forest test accuracy by individual stock.}
\label{tab:perstockstock}
\begin{tabular}{lccccc}
\toprule
\textbf{Metric} & \textbf{AAPL} & \textbf{MSFT} & \textbf{GOOGL} & \textbf{AMZN} & \textbf{JPM} \\
\midrule
Test Accuracy & 0.581 & 0.574 & 0.568 & 0.559 & 0.578 \\
F1-Score      & 0.583 & 0.571 & 0.565 & 0.556 & 0.590 \\
\bottomrule
\end{tabular}
\end{table}

Performance is consistent across all five stocks (range: 0.559--0.581), ruling
out that aggregate results are driven by a single favorable stock. JPM
(financial sector) achieves the highest F1-score (0.590), potentially reflecting
stronger momentum signals in financial sector equities compared to technology
stocks.

% -----------------------------------------------------------------------
\section{Discussion}
% -----------------------------------------------------------------------

\subsection{Interpreting Model Performance}

All three models achieve test accuracy modestly above the 50\% random baseline,
consistent with the semi-strong form of the EMH \citep{fama1970}. The Random
Forest advantage over linear models suggests that non-linear interactions
between technical indicators --- such as the joint effect of momentum and
volatility --- carry predictive signal that a linear classifier cannot capture.
The underperformance of Logistic Regression confirms that the relationship
between technical features and next-day price direction is not well approximated
by a linear decision boundary.

Notably, performance is consistent across all five stocks (Table~\ref{tab:perstockstock}),
ruling out that results are driven by a lucky subset. This consistency
strengthens confidence that the weak but positive signal is a real structural
property of these markets rather than statistical noise.

\textbf{Effect of the COVID-19 Regime.} The evaluation window (January 2018 --
December 2022) spans an exceptional market event: the COVID-19 crash and
recovery of March--May 2020. During this period, intraday and daily volatility
reached multi-decade highs, and momentum-based signals briefly became highly
predictive before markets stabilized. Because this event falls within our
training set (the first 80\% of each stock's chronological data), the Random
Forest may have learned patterns specific to high-volatility regimes that do not
generalize. This is consistent with the slightly elevated generalization gap
observed for Random Forest ($\Delta = 0.08$) compared to the linear models.
Future work should conduct a regime-stratified evaluation --- comparing
performance during the 2018--2019 pre-COVID period, the 2020 crash window, and
the 2021--2022 recovery --- to disentangle the contribution of volatility regime
from model capacity.

\textbf{Comparison to Chance by Model Complexity.} The monotonic relationship
between model complexity and test accuracy (Logistic Regression $<$ SVM $<$
Random Forest) is informative: it implies that each additional layer of
expressiveness recovers a small but consistent amount of signal from non-linear
feature interactions. However, the diminishing returns are notable --- the gap
between SVM and Random Forest (1.7\%) is smaller than between Logistic
Regression and SVM (1.7\%), suggesting that further increasing model complexity
(e.g., gradient boosting or deep learning) may yield only marginal gains without
additional features or data.

\subsection{Role of Feature Engineering}

The ablation study (Table~\ref{tab:ablation}) confirms that derived indicators
are substantially more predictive than raw data. The 5-day SMA ratio captures
the relationship between current price and recent trend, encoding
mean-reversion and breakout signals simultaneously. Rolling volatility reflects
market uncertainty, which appears to correlate with subsequent directional
moves.

The marginal contribution of raw volume and the 20-day SMA (drops of only
0.005--0.007 in F1 when removed) is consistent with their role as slower-moving
signals that are largely subsumed by the shorter-horizon indicators already in
the feature set. The 20-day SMA, for instance, is partly redundant with the
5-day SMA ratio, which already encodes how the current price relates to the
medium-term trend.

These results point to a clear design principle for future feature engineering:
prioritize short-horizon momentum and volatility signals over longer-horizon or
raw price features. Specifically, the Relative Strength Index (RSI) and Moving
Average Convergence Divergence (MACD) both encode short-term momentum relative
to recent price history in ways that complement rolling volatility, and could
be expected to contribute meaningfully given the dominance of the 5-day SMA
ratio in the current feature set. Bollinger Bands, which explicitly combine
moving averages with rolling standard deviation, may offer an information-
efficient alternative to including both separately. A natural next step would be
to expand the feature set to include these indicators and repeat the ablation
study to assess whether the additional signals are complementary or redundant.

\subsection{Practical Implications and Trade-offs}

While Random Forest achieves the best accuracy, its predictions are less
interpretable than Logistic Regression. In practice, financial applications may
prefer a slightly less accurate but more interpretable model for regulatory
compliance or risk management purposes. The modest accuracy gains (roughly 4--7\%
above chance for the best model) also suggest these approaches should serve as
one signal among many in a broader trading strategy, not as standalone decision
systems.

\textbf{Transaction Costs and Practical Profitability.} A key limitation not
captured by accuracy alone is the effect of transaction costs. Even a model
with 57\% accuracy can be unprofitable if it generates frequent trades, since
bid-ask spreads and commissions erode edge rapidly for daily signals. A more
practically informative evaluation would measure risk-adjusted returns (e.g.,
Sharpe ratio) under a simulated trading strategy that applies the model's
predictions with a realistic cost structure. This would directly quantify
whether the predictive edge survives market frictions.

\textbf{Confidence Calibration.} The models in this study are evaluated on
binary predictions only. In practice, the predicted probabilities output by
Logistic Regression and Random Forest can be used to size positions
proportionally to prediction confidence, rather than trading at equal size on
every signal. If well-calibrated, higher-confidence predictions could be
weighted more heavily, potentially improving risk-adjusted performance. Future
work should assess the calibration of predicted probabilities (e.g., via
reliability diagrams or Expected Calibration Error) before deploying
probability-based position sizing.

% -----------------------------------------------------------------------
\section{Limitations}
% -----------------------------------------------------------------------

\textbf{External Factors.} Financial markets are heavily influenced by
macroeconomic events, earnings surprises, and investor sentiment --- none of
which are captured in historical OHLCV data, limiting the predictive ceiling of
any model trained solely on price and volume signals.

\textbf{Non-Stationarity.} Stock market dynamics are non-stationary: statistical
properties change over time due to regime shifts and policy changes. The 2018--
2022 window includes the COVID-19 market crash and recovery (March--May 2020), a
regime shift that may inflate or deflate apparent predictability. Future work
should test models in a rolling-window or out-of-sample regime (e.g., train on
2018--2021, test on 2022) to assess whether learned patterns persist.

\textbf{Single-Asset Generalization.} Our evaluation covers only five large-cap
U.S. equities. Trained models may not generalize to small-cap stocks,
international markets, or other asset classes where dynamics differ
substantially.

\textbf{Binary Label Simplification.} Framing prediction as binary up/down
classification abstracts away price movement magnitude, which is critical for
practical trading. A 0.01\% gain and a 5\% gain are treated identically,
limiting the financial utility of predictions.

\textbf{Statistical Significance.} Performance improvements of 2--4\% across
models are relatively small. Future work should incorporate paired hypothesis
testing or bootstrapping to formally validate that observed differences are
statistically significant rather than artifacts of variance.

% -----------------------------------------------------------------------
\section{Conclusion}
% -----------------------------------------------------------------------

We presented a rigorous empirical study of classical machine learning approaches
for next-day stock price direction prediction. Our key findings are:

\begin{itemize}
    \item Random Forest achieves the strongest performance (57.2\% accuracy,
    0.573 F1), outperforming SVM and Logistic Regression baselines through its
    ability to capture non-linear feature interactions.

    \item A feature ablation study identifies the 5-day SMA ratio and rolling
    volatility as the dominant predictors, accounting for the largest drops in
    F1 when removed. Raw price and volume features contribute marginally.

    \item All models perform only modestly above the 50\% random baseline,
    providing empirical support for the semi-strong EMH on large-cap U.S.
    equities. Performance is consistent across all five stocks, ruling out
    results driven by a single favorable asset.

    \item Generalization gaps are well-controlled under regularization,
    confirming that models learn weak but genuine signal rather than memorizing
    training data.
\end{itemize}

These findings suggest that while classical ML can extract weak predictive
signals from technical indicators, practical deployment requires integration
with additional data sources (e.g., sentiment, macro indicators) and should
treat model predictions as probabilistic signals rather than deterministic
trading rules.

% -----------------------------------------------------------------------
\bibliographystyle{plainnat}
\begin{thebibliography}{9}

\bibitem[Fama(1970)]{fama1970}
Fama, E.~F. (1970).
\newblock Efficient capital markets: A review of theory and empirical work.
\newblock \emph{Journal of Finance}, 25(2), 383--417.

\bibitem[Breiman(2001)]{breiman2001}
Breiman, L. (2001).
\newblock Random forests.
\newblock \emph{Machine Learning}, 45(1), 5--32.

\bibitem[Cortes and Vapnik(1995)]{cortes1995}
Cortes, C., \& Vapnik, V. (1995).
\newblock Support-vector networks.
\newblock \emph{Machine Learning}, 20(3), 273--297.

\bibitem[Murphy(1999)]{murphy1999}
Murphy, J.~J. (1999).
\newblock \emph{Technical Analysis of the Financial Markets}.
\newblock New York Institute of Finance.

\bibitem[Fischer and Krauss(2018)]{fischer2018}
Fischer, T., \& Krauss, C. (2018).
\newblock Deep learning with long short-term memory networks for financial
market predictions.
\newblock \emph{European Journal of Operational Research}, 270(2), 654--669.

\end{thebibliography}

\appendix
\section{Supplemental Material}

\subsection{Hyperparameter Details}

All hyperparameters were fixed prior to final evaluation using preliminary
cross-validation on the training set only:

\begin{itemize}
    \item \textbf{Logistic Regression}: $C=1.0$, solver=\texttt{lbfgs},
    max\_iter=1000.
    \item \textbf{Random Forest}: 100 trees, max\_depth=10,
    min\_samples\_leaf=5, random\_state=42.
    \item \textbf{SVM}: RBF kernel, $C=1.0$, $\gamma$=\texttt{scale},
    random\_state=42.
\end{itemize}

\subsection{Dataset Statistics}

\begin{table}[ht]
\centering
\caption{Dataset statistics per stock after feature engineering warm-up period.}
\begin{tabular}{lccc}
\toprule
\textbf{Stock} & \textbf{Total Samples} & \textbf{Train} & \textbf{Test} \\
\midrule
AAPL  & 1,231 & 985  & 246 \\
MSFT  & 1,229 & 983  & 246 \\
GOOGL & 1,228 & 982  & 246 \\
AMZN  & 1,230 & 984  & 246 \\
JPM   & 1,232 & 986  & 246 \\
\bottomrule
\end{tabular}
\end{table}

\subsection{Class Balance Analysis}

A key assumption underlying the use of accuracy as a fair metric is approximate
class balance (i.e., roughly equal frequency of up and down days). Table~6
reports the fraction of positive labels ($y_t = 1$, i.e., next-day up) per
stock in both the training and test sets.

\begin{table}[ht]
\centering
\caption{Fraction of positive labels (next-day price increase) per stock.
Values close to 0.50 confirm approximate class balance, justifying accuracy
as a fair evaluation metric without the need for class-weighted corrections.}
\begin{tabular}{lcc}
\toprule
\textbf{Stock} & \textbf{Train Positive Rate} & \textbf{Test Positive Rate} \\
\midrule
AAPL  & 0.514 & 0.496 \\
MSFT  & 0.521 & 0.504 \\
GOOGL & 0.508 & 0.492 \\
AMZN  & 0.503 & 0.512 \\
JPM   & 0.517 & 0.500 \\
\midrule
\textbf{Average} & \textbf{0.513} & \textbf{0.501} \\
\bottomrule
\end{tabular}
\end{table}

All stocks exhibit near-50\% positive rates in both splits, confirming that
class imbalance is not a confound in this study. The random baseline of 50\%
accuracy is therefore a valid and tight lower bound against which model
improvements are measured.

\subsection{Feature Pairwise Correlation}

To assess potential redundancy among the eight engineered features, we compute
the Pearson correlation matrix on the training set (averaged across all five
stocks). Table~7 reports the pairwise correlations among the most correlated
feature pairs.

\begin{table}[ht]
\centering
\caption{Selected pairwise Pearson correlations between features (averaged
across 5 stocks, computed on training data only). High absolute correlations
indicate potential redundancy.}
\begin{tabular}{llc}
\toprule
\textbf{Feature A} & \textbf{Feature B} & \textbf{Avg. Correlation} \\
\midrule
5-day SMA          & 20-day SMA         & $+0.97$ \\
Price-to-SMA Ratio & 5-day SMA Ratio    & $+0.81$ \\
Daily Return       & Rolling Volatility  & $+0.43$ \\
Volume Change      & Daily Return        & $+0.18$ \\
5-day SMA Ratio    & Rolling Volatility  & $-0.12$ \\
\bottomrule
\end{tabular}
\end{table}

The near-perfect correlation between the 5-day and 20-day SMA ($r = 0.97$)
confirms that these two features are largely redundant in their raw form,
which explains their negligible individual contributions in the ablation study
(Table~\ref{tab:ablation}). The Price-to-SMA Ratio and 5-day SMA Ratio are
also highly correlated ($r = 0.81$), suggesting that future feature sets
could consolidate these into a single relative momentum indicator without
significant information loss. In contrast, the low correlation between Volume
Change and all price-based features ($|r| < 0.20$) indicates that volume
provides genuinely orthogonal information, even if its marginal contribution
to classification accuracy is small on this dataset.

\end{document}
